This Master's Thesis presents the design and implementation of a reproducible analytical platform for studying urban mobility patterns in New York City and forecasting demand for taxi and for-hire vehicle services. The primary source is the trip-record data published by the New York City Taxi and Limousine Commission (TLC), integrating Yellow Taxi, Green Taxi, For-Hire Vehicle (FHV), and High Volume For-Hire Vehicle (HVFHV) services. The analysis covers the 2018--2025 period, enabling a longitudinal comparison between pre-pandemic mobility, the disruption associated with COVID-19 in 2020, and the subsequent recovery.

The solution follows a Medallion Architecture. The Bronze layer preserves the original Parquet files, whereas the Silver and Gold layers are implemented in PostgreSQL and transformed with dbt. Data ingestion, scheduling and operational flow control are centralised in Apache NiFi, using process groups, parameters, explicit success and failure routes, retry policies and data-provenance capabilities. Docker Compose packages the required services and provides a reproducible deployment for NiFi, PostgreSQL, dbt, the Python machine-learning environment and Apache Superset. The Gold layer exposes dimensions, fact tables and analytical data marts intended to support demand analysis, mobility patterns, approximate congestion indicators, economic performance and driver-opportunity analysis.

Demand forecasting is formulated as a regression problem over the number of trips originating in each TLC zone and hourly time slot. Two forecasting horizons are considered: one hour and twenty-four hours. Three complementary supervised models are compared: Poisson regression, Random Forest and LightGBM\@. Evaluation is based on strictly temporal splits to avoid information leakage, using MAE, RMSE and WAPE\@. Final experimental figures are intentionally tied to the reproducible execution of the final repository rather than being estimated or manually introduced.
